\section{Notation and conventions}\label{app:notation}
\begin{center}\small
\begin{tabularx}{\linewidth}{@{}lX@{}}\toprule
Symbol & Meaning\\\midrule
\(N=3q,\Delta_N\)& Admitted ring length and native negative cycle Laplacian\\
\(\ell=3\lambda\)& Abbreviation for existing synchronizer strength\\
\(r,R,H_r,P,S_\phi,J_\phi\)& Positive radii, diagonal radii, symmetric pre-response, phase pre-response, synchronizer derivative, complete phase response\\
\(h,t,d\)& Triad scale-invariant algebraic quantities and \(d=1-3\ell\); t is not time\\
\(\Gamma_B\)& Bloch imaginary-coefficient obstruction\\
\(q,p= \re\Omega,\im\Omega\)& Real channel vectors; not ambient components\\
\(a,b,u,v\)& Their means and transverse components in the spatial-observable sections\\
\(N_f,T_f\)& Fixed normal/tangent frame matrices; not graph size or physical time\\
\(C_{\rm ch},W,\Gamma_{\rm ch}\)& Raw channel area and its horizontal-axial/vertical-polar projections\\
\(\mathcal H,S_G,A_G\)& Complex coherence, real Gram and antisymmetric pair-area matrix\\\bottomrule
\end{tabularx}
\end{center}
In the spectral sections a,b,c denote positive radii; in the observable sections a,b are channel means. The contexts and subscripts prevent an identification. The full mirror action is stated in \eqref{eq:reps}; the degree counts are not merely \(C_3\) counts.

The exact shell frames, with \(o=(0,\sqrt3/6,0)\), are
\[
\begin{array}{c|c|c|c}
i&c_i&n_i&t_i\\\hline
A&(-1/4,\sqrt3/4,0)&(-\sqrt3/2,1/2,0)&(-1/2,-\sqrt3/2,0)\\
B&(0,0,0)&(0,-1,0)&(1,0,0)\\
C&(1/4,\sqrt3/4,0)&(\sqrt3/2,1/2,0)&(-1/2,\sqrt3/2,0)
\end{array}
\]
Direct dot products prove \(N_f^TN_f=T_f^TT_f=(3/2)\Pi\). At a general real fold angle beta, the accepted panel maps have material derivatives
\[
\partial_uX_1=(-\cos\beta,-\sin\beta,0),\quad
\partial_uX_2=(1,0,0),\quad
\partial_uX_3=(-\cos\beta,\sin\beta,0),\qquad
\partial_zX_i=e_z.
\]
Each horizontal derivative has unit norm and is orthogonal to \(e_z\); their pullback metrics are therefore \(I_2\). This local identity does not assert an embedded closed assembly at arbitrary fold angle. Embedding folds do not imply a variable intrinsic panel metric.

For completeness, the preparation Fourier vectors are
\(f_0=e/\sqrt3,\ f_1=(1,\bar\omega,\omega)^T/\sqrt3,\ f_2=\bar f_1\),
\(\omega=e^{2\pi i/3}\). With \(p_0=f_0f_0^\dagger\),
\[
A_s=[q_sI+(r_s-q_s)p_0]\diag(1,\bar\omega,\omega)
[c_sI+(1-c_s)p_0].
\]
The fixed two-factor matrix is
\[
B_s=\diag(e^{-i\alpha},e^{i\alpha})
\begin{pmatrix}\cos\beta&-i\sin\beta\\-i\sin\beta&\cos\beta\end{pmatrix},
\quad\alpha=2\pi(.244),\quad\beta=\pi(.244).
\]
This fully specifies the algebra used in the type chain, without assigning material boundary meaning.

\section{Complete theorem and claim ledger}\label{app:ledger}
This is the anti-overstatement ledger. Predicate ordinals resolve to the named lists in \texttt{verification/retained\_predicates.json}; checks do not replace proof locations.\par
\begingroup\footnotesize\raggedright
\par\noindent\begin{minipage}{\linewidth}
\subsection*{L01 Native map and regular chart}
\textbf{DEFINITION / ADOPTION}. The native map is amplitude/coupling followed by harmonic-three phase synchronization; Arg0(0)=0. N=3q, coefficients repeated every three nodes.\par
\textbf{Domain:} Real parameters; finite states; differential claims only on nonzero pre-sync charts.\par
\textbf{Authority/proof:} GR0 \S{}5.1--5.2; native inventory \S{}3; main text \S\ref{sec:native}.\par
\textbf{Prior ownership:} Yes: Paper A \S{}6.1\par
\textbf{Qualification:} No global differentiability at zeros for nonzero lambda.\par
\textbf{Interpretation limit:} No physical time or spatial mesh is introduced.\par
\textbf{Supporting predicates:} GR0 ordinals 33--35\par
\end{minipage}\par\medskip
\par\noindent\begin{minipage}{\linewidth}
\subsection*{L02 Panel and state isometries}
\textbf{DERIVED IDENTITY}. The panel pullback metric is \(I_2\). The decoder \(D_i(q+ip)=qt_i+pe_z\) is isometric; \(ED=I\), while \(DE\) is the ambient tangent projector.\par
\textbf{Domain:} Canonical frames; arbitrary real fold for individual panel isometry; D/E restricted to tangent domain.\par
\textbf{Authority/proof:} GR0 \S{}4(a); main text \S\ref{sec:attachment}.\par
\textbf{Prior ownership:} Yes: Papers B/C; Atlas 03/04\par
\textbf{Qualification:} DE is not the ambient nine-dimensional identity.\par
\textbf{Interpretation limit:} Fixed panel geometry and state geometry remain distinct.\par
\textbf{Supporting predicates:} GR0 ordinals 1, 3, 5, 7--8, 12--13, 17--18\par
\end{minipage}\par\medskip
\par\noindent\begin{minipage}{\linewidth}
\subsection*{L03 Matched transport}
\textbf{EXACT THEOREM}. \(T_{ij}=t_it_j^T+e_ze_z^T\) satisfies \(T_{ij}T_{jk}=T_{ik}\). Closed products restrict to identity on \(W_i\).\par
\textbf{Domain:} Labelled tangent fibres; ambient maps have rank two.\par
\textbf{Authority/proof:} GR0 \S{}4(b)--(c); SA0 \S{}3.2; main text \S\ref{sec:attachment}.\par
\textbf{Prior ownership:} Yes: Paper B Proposition 2\par
\textbf{Qualification:} Ambient closed product is a projector.\par
\textbf{Interpretation limit:} No arbitrary surface-path or physical connection is derived.\par
\textbf{Supporting predicates:} GR0 ordinals 9--11, 14--16, 19--22\par
\end{minipage}\par\medskip
\par\noindent\begin{minipage}{\linewidth}
\subsection*{L04 Potential and exact intensity budget}
\textbf{DERIVED IDENTITY}. The pre-stage increment is \(-\nabla\mathcal V\), with \(\mathcal V=\eps\sum_i(s_i^2/4-k_is_i/2)+(g/2)\sum_{\{i,j\}}|z_i-z_j|^2\), \(s_i=|z_i|^2\). The complete intensity increment includes the squared norm of the pre-increment.\par
\textbf{Domain:} Triad; real Cartesian gradient; exact normalization.\par
\textbf{Authority/proof:} GR0 \S{}4(d)--(e); main text \S\ref{sec:native}.\par
\textbf{Prior ownership:} Yes: Papers B/D/E; Atlas 07\par
\textbf{Qualification:} No universal discrete descent: common input 2 at eps=k=1,g=0 maps to -4.\par
\textbf{Interpretation limit:} Not a calibrated physical energy or elastic shell energy.\par
\textbf{Supporting predicates:} GR0 ordinals 23--29, 31--32\par
\end{minipage}\par\medskip
\par\noindent\begin{minipage}{\linewidth}
\subsection*{L05 Positive synchronized fixed circle}
\textbf{EXACT THEOREM}. For \(\eps\ne0\), fixed states on the nonzero common-channel line form \(\sqrt{k}e^{i\chi}\one\) exactly when \(k>0\).\par
\textbf{Domain:} Any admitted finite ring; real g,lambda.\par
\textbf{Authority/proof:} GR0 \S{}5.1; main text \S\ref{sec:homogeneous}.\par
\textbf{Prior ownership:} Yes in unit/triad cases; general equal-k formulation consolidated by GR0\par
\textbf{Qualification:} At epsilon=0 every common amplitude is fixed and an extra neutral direction appears.\par
\textbf{Interpretation limit:} A fixed circle in state space is not a material orbit.\par
\textbf{Supporting predicates:} Analytic proof/source statement; no individually matched predicate claimed.\par
\end{minipage}\par\medskip
\par\noindent\begin{minipage}{\linewidth}
\subsection*{L06 Complete homogeneous derivative}
\textbf{DERIVED IDENTITY}. The amplitude block is \((1-2\eps k)I+g\Delta\); the phase block is \((I+3\lambda\Delta)(I+g\Delta)\).\par
\textbf{Domain:} Positive equal-k background; local real amplitude/phase chart.\par
\textbf{Authority/proof:} GR0 \S{}5.2; main text \S\ref{sec:homogeneous}.\par
\textbf{Prior ownership:} Unit triad known in Paper F; all-ring equal-k derivation GR0\par
\textbf{Qualification:} Stage order matters away from homogeneous commutation.\par
\textbf{Interpretation limit:} No inertia is added.\par
\textbf{Supporting predicates:} GR0 ordinals 33--35, 45\par
\end{minipage}\par\medskip
\par\noindent\begin{minipage}{\linewidth}
\subsection*{L07 Finite Fourier spectrum}
\textbf{EXACT THEOREM}. \(\mu_a=1-2\eps k-gh_m\), \(\mu_\phi=(1-gh_m)(1-3\lambda h_m)\), and \(h_m=4\sin^2(\pi m/N)\).\par
\textbf{Domain:} N=3q, m=0,...,N-1, equal positive k.\par
\textbf{Authority/proof:} GR0 \S{}5.3; main text \S\ref{sec:homogeneous}.\par
\textbf{Prior ownership:} Cycle Fourier calculus and unit triad known; present synthesis from GR0\par
\textbf{Qualification:} Triad pullback modes have h=0,3,3 and are not long-wave modes.\par
\textbf{Interpretation limit:} Real/even multipliers exclude propagating linear wave dispersion only on this branch.\par
\textbf{Supporting predicates:} GR0 ordinals 36--41\par
\end{minipage}\par\medskip
\par\noindent\begin{minipage}{\linewidth}
\subsection*{L08 Normal stability}
\textbf{EXACT THEOREM}. Strict normal contraction is equivalent to \(|\mu_a(m)|<1\) for every m and \(|\mu_\phi(m)|<1\) for nonzero m; it yields local finite-N attraction modulo common phase.\par
\textbf{Domain:} Fixed finite N; nonzero chart; strict inequalities.\par
\textbf{Authority/proof:} GR0 \S{}5.4; main text \S\ref{sec:homogeneous}.\par
\textbf{Prior ownership:} Prior local methods in Paper F; scope GR0\par
\textbf{Qualification:} No uniform-in-N basin; equality at unit modulus is undecided by the linear test.\par
\textbf{Interpretation limit:} Relaxation may include sign alternation.\par
\textbf{Supporting predicates:} GR0 ordinals 36--41\par
\end{minipage}\par\medskip
\par\noindent\begin{minipage}{\linewidth}
\subsection*{L09 Restricted modal limit}
\textbf{EXACT THEOREM}. For fixed m and \(D_\phi=g+3\lambda>0\), \(\mu_\phi(2\pi m/N)^{\lfloor\tau N^2\rfloor}\) tends uniformly on bounded nonnegative tau intervals to \(e^{-4\pi^2D_\phi m^2\tau}\), with \(O(N^{-2})\) error.\par
\textbf{Domain:} Fixed native coefficients and fixed finite collection of Fourier indices; N$\to$infinity through 3q.\par
\textbf{Authority/proof:} GR0 \S{}6.1--6.2; main text \S\ref{sec:homogeneous}.\par
\textbf{Prior ownership:} Not asserted as prior result; GR0\par
\textbf{Qualification:} No nonlinear continuum theorem or control of all growing modes; D\_phi=0 has quartic leading decrement.\par
\textbf{Interpretation limit:} Diffusion-like modal analogy, no physical dt or length.\par
\textbf{Supporting predicates:} GR0 ordinals 42--44\par
\end{minipage}\par\medskip
\par\noindent\begin{minipage}{\linewidth}
\subsection*{L10 Positive unequal branch and screening}
\textbf{EXACT THEOREM}. The local analytic positive solution of \(\eps r_i(k_i-r_i^2)+g(\Delta r)_i=0\) has \((2\eps kI-g\Delta)u=\eps\sqrt{k}\,\kappa\). For three-periodic zero-mean kappa, \(u=\eps\sqrt{k}\,\kappa/(2\eps k+3g)\).\par
\textbf{Domain:} k>0; strict finite-N normal stability at equal point; k\_i=k+eta kappa\_i.\par
\textbf{Authority/proof:} GR1 \S{}2, (2)--(6); main text \S\ref{sec:unequal}.\par
\textbf{Prior ownership:} Fixed equation known Paper A \S{}7.5; response formula GR1\par
\textbf{Qualification:} Positive inverse derivative, local uniqueness only; no global branch continuation.\par
\textbf{Interpretation limit:} Screened-Poisson structural analogy only when g>0.\par
\textbf{Supporting predicates:} GR1 ordinals 1, 5--6\par
\end{minipage}\par\medskip
\par\noindent\begin{minipage}{\linewidth}
\subsection*{L11 Pre-stage weighted response}
\textbf{EXACT THEOREM}. \(P=R^{-1}H_rR\), \(H_rr=r\), and \(R^2P=P^TR^2\). The inner product has weights \(r_i^2\) and Dirichlet form \(g\sum_{\{i,j\}}r_ir_j(f_i-f_j)^2\).\par
\textbf{Domain:} Every positive fixed branch; real g; uniqueness of weights up to scale when g nonzero and graph connected.\par
\textbf{Authority/proof:} GR1 \S{}3, (7)--(12); main text \S\ref{sec:weighted}.\par
\textbf{Prior ownership:} Derivative framework prior; explicit weighting/classification GR1\par
\textbf{Qualification:} g=0 permits every positive diagonal weight; g<0 reverses form sign; stochasticity needs entrywise nonnegativity.\par
\textbf{Interpretation limit:} Background-dependent mathematical response geometry, not spacetime.\par
\textbf{Supporting predicates:} GR1 ordinals 1--4\par
\end{minipage}\par\medskip
\par\noindent\begin{minipage}{\linewidth}
\subsection*{L12 Full unequal derivative}
\textbf{DERIVED IDENTITY}. \(J_\phi=S_\phi P\), with \(S_\phi=I+\ell\Delta\), \(\ell=3\lambda\); the amplitude block is \(H_r-2\eps R^2\). There is no first-order amplitude/phase mixing.\par
\textbf{Domain:} Positive branch; phase coordinates b=R phi.\par
\textbf{Authority/proof:} GR1 \S{}3, (13)--(15); main text \S\ref{sec:weighted}.\par
\textbf{Prior ownership:} Prior derivative framework; GR1 full formulation\par
\textbf{Qualification:} Cartesian imaginary block is R S\_phi R\textasciicircum{}-1 H; not SP in those coordinates.\par
\textbf{Interpretation limit:} Two separately self-adjoint factors need not have a common metric.\par
\textbf{Supporting predicates:} GR1 ordinals 7--12\par
\end{minipage}\par\medskip
\par\noindent\begin{minipage}{\linewidth}
\subsection*{L13 Triad cycle factorization}
\textbf{DERIVED IDENTITY}. \(\mathcal C_{\rm cyc}=g\ell(1-3\ell)(a-b)(a-c)(b-c)\mathcal K/(a^2b^2c^2)\), where \(\mathcal K=g^2\ell\sigma_1^3+[g^2(1-2\ell)-g\ell]\sigma_1\sigma_2+(3g\ell-g+\ell)\sigma_3\).\par
\textbf{Domain:} a,b,c>0; exact algebra for any real g,ell.\par
\textbf{Authority/proof:} GR1 \S{}4, (16)--(20); main text \S\ref{sec:cycle}.\par
\textbf{Prior ownership:} GR1; no historical priority claim\par
\textbf{Qualification:} GR2 (36) omits a '+' in its restatement; this paper uses verified GR1 (20).\par
\textbf{Interpretation limit:} Algebraic response circulation, not physical current.\par
\textbf{Supporting predicates:} GR1 ordinals 13\par
\end{minipage}\par\medskip
\par\noindent\begin{minipage}{\linewidth}
\subsection*{L14 Triad diagonal criterion and exceptions}
\textbf{EXACT THEOREM}. With positive off-diagonals, \(\mathcal C_{\rm cyc}=0\) iff positive diagonal balance exists. One choice is \(W=\diag(J_{21}J_{31},J_{12}J_{31},J_{13}J_{21})\). Exceptional factors are \(g=0\), \(\ell=0\), \(\ell=1/3\), equal radius pairs, or \(\mathcal K=0\).\par
\textbf{Domain:} Sufficiently near a strictly stable equal triad, guaranteeing positive off-diagonals.\par
\textbf{Authority/proof:} GR1 \S{}4, (18)--(24); main text \S\ref{sec:cycle}.\par
\textbf{Prior ownership:} GR1 classification; cycle logic standard and reproved\par
\textbf{Qualification:} Outside positive-entry neighbourhood zero cycle alone does not prove positive weights.\par
\textbf{Interpretation limit:} Locality of weights is a mathematical restriction.\par
\textbf{Supporting predicates:} GR1 ordinals 14--27\par
\end{minipage}\par\medskip
\par\noindent\begin{minipage}{\linewidth}
\subsection*{L15 Cubic onset}
\textbf{DERIVED IDENTITY}. \(\mathcal C_{\rm cyc}=g\ell(1-3\ell)K_0[\eps/(2\eps k+3g)]^3V(\kappa)\eta^3+O(\eta^4)\), with \(K_0=g(9g-1)+\ell(1-3g)^2\).\par
\textbf{Domain:} Weak branch of result 10; three-periodic zero-mean perturbation.\par
\textbf{Authority/proof:} GR1 \S{}5, (25)--(26); main text \S\ref{sec:cycle}.\par
\textbf{Prior ownership:} GR1\par
\textbf{Qualification:} If K0=0 the obstruction is O(eta\textasciicircum{}5) or smaller, not automatically zero.\par
\textbf{Interpretation limit:} First-order fitted weights cannot settle exact diagonal balance.\par
\textbf{Supporting predicates:} GR1 ordinals 13, 15\par
\end{minipage}\par\medskip
\par\noindent\begin{minipage}{\linewidth}
\subsection*{L16 Larger-ring diagonal impossibility}
\textbf{EXACT THEOREM}. For \(N\ge6\) and \(g\ell\ne0\), distinct positive repeated radii forbid positive diagonal balance. For \((a,a,c)\), it exists iff \(\ell-g-g\ell(2a/c+c/a)=0\), and then \(W=R\).\par
\textbf{Domain:} N=3q>=6; positive three-periodic branch; no entrywise positivity/stability assumption.\par
\textbf{Authority/proof:} GR1 \S{}6, (27)--(30); main text \S\ref{sec:cycle}.\par
\textbf{Prior ownership:} GR1\par
\textbf{Qualification:} g=0,ell=0,or equal radii remain reversible; triad exceptional loci do not suffice on the full ring.\par
\textbf{Interpretation limit:} Distance-two response links strengthen the node-local obstruction.\par
\textbf{Supporting predicates:} GR1 ordinals 28--38\par
\end{minipage}\par\medskip
\par\noindent\begin{minipage}{\linewidth}
\subsection*{L17 SPD criterion}
\textbf{EXACT THEOREM}. A real J admits SPD G with \(GJ=J^TG\) iff its spectrum is real and semisimple. All solutions are \(G=V^{-T}KV^{-1}\), with K SPD and commuting with the real eigenvalue matrix. Positive sums of \(E_\rho^TE_\rho\) also construct solutions.\par
\textbf{Domain:} Any finite real matrix; repeated semisimple/zero roots allowed.\par
\textbf{Authority/proof:} GR2 \S{}2, (4)--(6); main text \S\ref{sec:spd}.\par
\textbf{Prior ownership:} Standard finite-dimensional lemma, reproved; not a novelty claim\par
\textbf{Qualification:} Defective real matrices fail as well as complex-spectrum matrices.\par
\textbf{Interpretation limit:} G is nonunique and need not be node-local or physically privileged.\par
\textbf{Supporting predicates:} Analytic proof/source statement; no individually matched predicate claimed.\par
\end{minipage}\par\medskip
\par\noindent\begin{minipage}{\linewidth}
\subsection*{L18 Exact triad reality}
\textbf{EXACT THEOREM}. \(M_3=\sigma R^{-1}-\one\one^T\) is symmetric positive semidefinite with kernel r; its positive roots have sum h and product t. The quotient polynomial is \(Q(\xi)=\xi^2-d(2-gh)\xi+d^2(1-gh+g^2t)\), with \(\mathcal D_3=d^2g^2(h^2-4t)\ge0\).\par
\textbf{Domain:} Positive a,b,c; arbitrary real g,ell; d=1-3ell; transverse means quotient by common phase.\par
\textbf{Authority/proof:} GR2 \S{}3, (7)--(13); main text \S\ref{sec:triad}.\par
\textbf{Prior ownership:} GR2\par
\textbf{Qualification:} h\textasciicircum{}2-4t=0 iff radii equal; g=0 or d=0 also gives semisimple transverse coincidence.\par
\textbf{Interpretation limit:} Triad diagonal obstruction does not force complex roots.\par
\textbf{Supporting predicates:} GR2 ordinals 1--5\par
\end{minipage}\par\medskip
\par\noindent\begin{minipage}{\linewidth}
\subsection*{L19 Full triad resonance defect}
\textbf{EXACT THEOREM}. For \(M_3w_\alpha=\mu_\alpha w_\alpha\), \(w_\alpha\perp r\), set \(v_\alpha=R^{-1}w_\alpha\), \(p_\alpha=1-g\mu_\alpha\). A full defect occurs at \(dp_\alpha=1\) with \(\ell p_\alpha\one^Tv_\alpha\ne0\).\par
\textbf{Domain:} Positive nonconstant background; degenerate g,ell,d cases treated separately.\par
\textbf{Authority/proof:} GR2 \S{}3, (14)--(15); main text \S\ref{sec:triad}.\par
\textbf{Prior ownership:} GR2\par
\textbf{Qualification:} Three distinct radii imply a defective resonance; equal-pair antisymmetric collision is semisimple. Local strict stability excludes eigenvalue-one resonance.\par
\textbf{Interpretation limit:} Transverse quotient semisimplicity is not always full-matrix semisimplicity.\par
\textbf{Supporting predicates:} GR2 ordinals 17--18\par
\end{minipage}\par\medskip
\par\noindent\begin{minipage}{\linewidth}
\subsection*{L20 Triad discriminant onset}
\textbf{DERIVED IDENTITY}. \(\mathcal D_3=6d^2g^2[\eps/(2\eps k+3g)]^2\sum_i\kappa_i^2\eta^2+O(\eta^3)\). No order produces complex transverse roots on a positive triad.\par
\textbf{Domain:} Weak branch; dg and kappa nonzero for positive leading term.\par
\textbf{Authority/proof:} GR2 \S{}4, (16); main text \S\ref{sec:triad}.\par
\textbf{Prior ownership:} GR2\par
\textbf{Qualification:} dg=0 leaves exact real coincidence.\par
\textbf{Interpretation limit:} Independent from the cubic diagonal-cycle obstruction.\par
\textbf{Supporting predicates:} GR2 ordinals 12\par
\end{minipage}\par\medskip
\par\noindent\begin{minipage}{\linewidth}
\subsection*{L21 Finite Bloch blocks}
\textbf{DERIVED IDENTITY}. For \(\phi_{n,\alpha}=z^nv_\alpha\), \(z^q=1\), the exact block is \(J(z)=(I+\ell D(z))R^{-1}H_r(z)R\), where the corner entries of the cycle block \(D(z)\) are \(z^{-1},z\), its other off-diagonals are one and its diagonal is minus two.\par
\textbf{Domain:} N=3q; repeated a,b,c>0; exact finite Fourier transform under translation by three.\par
\textbf{Authority/proof:} GR2 \S{}5, (17); main text \S\ref{sec:bloch}.\par
\textbf{Prior ownership:} GR2 specialization; finite Fourier decomposition standard\par
\textbf{Qualification:} Conjugate blocks pair; each block need not have real coefficients.\par
\textbf{Interpretation limit:} No continuum approximation.\par
\textbf{Supporting predicates:} GR2 ordinals 19, 21, 23--24, 26, 28, 30, 32\par
\end{minipage}\par\medskip
\par\noindent\begin{minipage}{\linewidth}
\subsection*{L22 Bloch polynomial and circulation coefficient}
\textbf{DERIVED IDENTITY}. \(\det(\xi I-J(z))=\xi^3-T_B\xi^2+B_B(z)\xi-D_B(z)\), with coefficients in \eqref{eq:BT}--\eqref{eq:BD}. In particular \(\Gamma_B=\ell g[\ell(1+3g)-g](a-b)(a-c)(b-c)/(abc)\) and \(\im B_B=-\Gamma_B\sin\theta\).\par
\textbf{Domain:} Same finite positive branch; all real parameters; |z|=1.\par
\textbf{Authority/proof:} GR2 \S{}5, (18)--(20); main text \S\ref{sec:bloch}.\par
\textbf{Prior ownership:} GR2\par
\textbf{Qualification:} Gamma\_B differs from the triad cycle factor and the chiral scalar Gamma\_ch.\par
\textbf{Interpretation limit:} Spectral response coefficient, not physical flux.\par
\textbf{Supporting predicates:} GR2 ordinals 6--11\par
\end{minipage}\par\medskip
\par\noindent\begin{minipage}{\linewidth}
\subsection*{L23 Exact block decision}
\textbf{EXACT THEOREM}. Nonzero \(\Gamma_B\sin\theta\) forces non-real roots. For real cubic coefficients, positive discriminant gives three simple real roots, negative gives a complex pair, and zero requires a rank/minimal-polynomial test.\par
\textbf{Domain:} Every admitted Bloch block; full ring SPD iff all blocks have real semisimple spectra.\par
\textbf{Authority/proof:} GR2 \S{}5, (21)--(23); main text \S\ref{sec:bloch}.\par
\textbf{Prior ownership:} GR2 application; cubic/rank tests reproved\par
\textbf{Qualification:} At a double root rho with distinct sigma, rank(J-rho I)=1 is required; at triple root J=rho I is required.\par
\textbf{Interpretation limit:} Non-real does not imply defective or unstable.\par
\textbf{Supporting predicates:} GR2 ordinals 20, 22, 25, 27, 29, 31\par
\end{minipage}\par\medskip
\par\noindent\begin{minipage}{\linewidth}
\subsection*{L24 Generic larger-ring complex response}
\textbf{EXACT THEOREM}. For \(q\ge3\), three distinct radii and \(g\ell[\ell(1+3g)-g]\ne0\) force non-real spectrum. \(\Gamma_B\) is cubic in weak inequality; simple nonzero roots have the cubic imaginary shift in Section~\ref{sec:bloch}.\par
\textbf{Domain:} N>=9; repeated weak branch; simple-root expansion needs distinct nonzero equal-background root.\par
\textbf{Authority/proof:} GR2 \S{}5, (24)--(27); main text \S\ref{sec:bloch}.\par
\textbf{Prior ownership:} GR2\par
\textbf{Qualification:} Vanishing Gamma\_B is not sufficient for real semisimple spectrum; stability can persist.\par
\textbf{Interpretation limit:} Oscillatory linear response only.\par
\textbf{Supporting predicates:} GR2 ordinals 9, 27, 31\par
\end{minipage}\par\medskip
\par\noindent\begin{minipage}{\linewidth}
\subsection*{L25 Generic weak six-ring reality}
\textbf{DERIVED IDENTITY}. At \(z=-1\), \(\mathfrak D_\pi=6g^2(1-\ell)^2\sum_i v_i^2(3E_\pi)^4\eta^2+O(\eta^3)\), where \(E_\pi=g+\ell-5g\ell\).\par
\textbf{Domain:} Fixed stable equal N=6 baseline; g(1-ell)Epi nonzero and nontrivial inequality.\par
\textbf{Authority/proof:} GR2 \S{}6, (29)--(30); main text \S\ref{sec:six}.\par
\textbf{Prior ownership:} GR2\par
\textbf{Qualification:} Neighbourhood not uniform near exceptional baselines. At Epi=0 a separate first-order test applies; ell=1 needs rank test.\par
\textbf{Interpretation limit:} Six and twelve rings need not share spectral reality.\par
\textbf{Supporting predicates:} GR2 ordinals 13--16\par
\end{minipage}\par\medskip
\par\noindent\begin{minipage}{\linewidth}
\subsection*{L26 Exact six-ring nilpotent fixture}
\textbf{EXACT THEOREM}. For \(r=(1-\tau,1,1+\tau)\), \(g=(1-\tau^2)/4\), \(\ell=1\), the exact block is \(J(-1)=(\tau^2/4)(1,-1,1)^T(1+\tau,2,1-\tau)\). It has rank one, square zero, and Jordan sizes \(2+1\).\par
\textbf{Domain:} 0<|tau| sufficiently small; eps=1/20; k\_i from exact native fixed equations.\par
\textbf{Authority/proof:} GR2 \S{}6, (31)--(33); main text \S\ref{sec:six}.\par
\textbf{Prior ownership:} GR2\par
\textbf{Qualification:} At tau=1/1000 the exact contraction box links fixture to weak native branch; varying tau does not hold kappa fixed.\par
\textbf{Interpretation limit:} A normally stable response can be defective.\par
\textbf{Supporting predicates:} GR2 ordinals 33--34, 36, 41--42\par
\end{minipage}\par\medskip
\par\noindent\begin{minipage}{\linewidth}
\subsection*{L27 Six-ring complex unfolding}
\textbf{EXACT THEOREM}. With the same \(r,g,k_i,\eps\) and \(\ell=1+\delta\), the coefficients are \eqref{eq:unfoldcoeff}, and \(\mathfrak D_\pi=(27/2)\tau^{12}\delta^3+O(\delta^4)\).\par
\textbf{Domain:} Fixed nonzero small tau; sufficiently small delta.\par
\textbf{Authority/proof:} GR2 \S{}6, (34)--(35); main text \S\ref{sec:six}.\par
\textbf{Prior ownership:} GR2\par
\textbf{Qualification:} delta<0 complex pair; delta=0 defect; delta>0 distinct real roots; exact rational signs checked at $\pm$10\textasciicircum{}-14 for tau=10\textasciicircum{}-3.\par
\textbf{Interpretation limit:} Modes can remain inside unit disk.\par
\textbf{Supporting predicates:} GR2 ordinals 35, 37--38\par
\end{minipage}\par\medskip
\par\noindent\begin{minipage}{\linewidth}
\subsection*{L28 Actual spatial representation}
\textbf{DEFINITION / ADOPTION}. The generators act by \(C_3:(q,p)\mapsto(Pq,Pp)\), \(\sigma_v:(q,p)\mapsto(-Q_vq,Q_vp)\), and \(\sigma_h:(q,p)\mapsto(q,-p)\).\par
\textbf{Domain:} Fixed accepted shell frames; spatial polar/axial laws distinct.\par
\textbf{Authority/proof:} CM0 \S{}2, actual mirror-action table; main text \S\ref{sec:observables}.\par
\textbf{Prior ownership:} Yes: Paper B conventions, reused by CM0\par
\textbf{Qualification:} Not ordinary S3 action on both channel triples. This is kinematic covariance; F carries k under permutations and has a regular-domain qualification for mirrors involving common phase.\par
\textbf{Interpretation limit:} Kinematic symmetry does not silently give full fixed-unequal-k dynamical symmetry.\par
\textbf{Supporting predicates:} CM0 ordinals 1--39\par
\end{minipage}\par\medskip
\par\noindent\begin{minipage}{\linewidth}
\subsection*{L29 Ambient polynomial multiplicities}
\textbf{EXACT THEOREM}. Real-linear polar and axial map dimensions are 2 and 2; homogeneous quadratic dimensions are 5 and 3. Complete bases appear in Theorem~\ref{thm:observables}.\par
\textbf{Domain:} Real polynomials in six state coordinates, full D3h, fixed geometry.\par
\textbf{Authority/proof:} CM0 \S{}3; \S{}4 for sectors; main text \S\ref{sec:observables}.\par
\textbf{Prior ownership:} CM0 census; no priority claim\par
\textbf{Qualification:} C3 alone admits many more maps; mixtures' images are not sums of separate images.\par
\textbf{Interpretation limit:} No unique core observable follows.\par
\textbf{Supporting predicates:} CM0 ordinals 4--44\par
\end{minipage}\par\medskip
\par\noindent\begin{minipage}{\linewidth}
\subsection*{L30 Optional common-phase invariance}
\textbf{EXACT THEOREM}. The optional common-phase-invariant quadratic polar space is spanned by \(N_f(au+bv)\), \(N_f(u^2+v^2)\), \(\Gamma_{\rm ch}e_z\). The axial invariant space is spanned by W.\par
\textbf{Domain:} Additional common-phase invariance imposed as a kinematic requirement.\par
\textbf{Authority/proof:} CM0 \S{}3, (8)--(9); main text \S\ref{sec:observables}.\par
\textbf{Prior ownership:} CM0 classification; W definition already accepted\par
\textbf{Qualification:} There is no nonzero invariant real-linear map; this is not a globally valid phase quotient of F at zeros.\par
\textbf{Interpretation limit:} Even extra phase invariance does not select a polar position observable.\par
\textbf{Supporting predicates:} CM0 ordinals 45\par
\end{minipage}\par\medskip
\par\noindent\begin{minipage}{\linewidth}
\subsection*{L31 Independent chiral projections}
\textbf{EXACT THEOREM}. \(C_{\rm ch}=q\times p=e\times(av-bu)+u\times v\). Its horizontal-axial projection is \(W=T_fC_{\rm ch}=\sqrt3(bN_fu-aN_fv)\); its vertical-polar projection is \((e^TC_{\rm ch})e_z=\Gamma_{\rm ch}e_z\).\par
\textbf{Domain:} D3h shell representation; degree-two antisymmetric bilinears.\par
\textbf{Authority/proof:} CM0 \S{}4, (13); SA0 \S{}6; main text \S\ref{sec:observables}.\par
\textbf{Prior ownership:} C and W accepted earlier; complete two-projection statement CM0\par
\textbf{Qualification:} Neither projection alone determines C; raw C is not ambient; axial parity need not imply chirality.\par
\textbf{Interpretation limit:} No physical rotation axis.\par
\textbf{Supporting predicates:} CM0 ordinals 46--47\par
\end{minipage}\par\medskip
\par\noindent\begin{minipage}{\linewidth}
\subsection*{L32 Observable kernels and images}
\textbf{DERIVED IDENTITY}. For transverse u,v, \(\|N_f(u^2)\|^2=\|u\|^4/4\), \(\|N_f(uv)\|^2=\|u\|^2\|v\|^2/4\), and \(\Gamma_{\rm ch}^2=3(\|u\|^2\|v\|^2-(u\cdot v)^2)\). Products inside \(N_f\) are entrywise.\par
\textbf{Domain:} u,v perpendicular to e; native fixed N.\par
\textbf{Authority/proof:} CM0 \S{}4, (10)--(12) and generator tables; main text \S\ref{sec:observables}.\par
\textbf{Prior ownership:} CM0\par
\textbf{Qualification:} Nonlinear kernels are zero sets; mixture images must be computed jointly.\par
\textbf{Interpretation limit:} Nonzero spatial readout at a fixed state is not displacement per step.\par
\textbf{Supporting predicates:} CM0 ordinals 48--50\par
\end{minipage}\par\medskip
\par\noindent\begin{minipage}{\linewidth}
\subsection*{L33 Observable increments and response}
\textbf{DERIVED IDENTITY}. \(\Delta V=V(F(\Omega))-V(\Omega)\). Polynomial increments follow substitution; local response at a fixed background is \(DV(\Omega_*)DF(\Omega_*)\).\par
\textbf{Domain:} Degree<=2 basis; equal or weak unequal positive branch for linear response.\par
\textbf{Authority/proof:} CM0 \S{}5, (14)--(20) and response tables; main text \S\ref{sec:observables}.\par
\textbf{Prior ownership:} CM0 application of chain rule\par
\textbf{Qualification:} Generic basis readouts can be generated from zero; not preserved or autonomous in general.\par
\textbf{Interpretation limit:} Delta V is not velocity or acceleration.\par
\textbf{Supporting predicates:} CM0 ordinals 53--78\par
\end{minipage}\par\medskip
\par\noindent\begin{minipage}{\linewidth}
\subsection*{L34 Invertible shell attachment}
\textbf{EXACT THEOREM}. \(D:\C^3\to\bigoplus_i W_i\) has inverse E. The full section evolves by \(DFE\), with transported graph increment and outward-normal phase rotation.\par
\textbf{Domain:} Exactly three labelled tangent fibres; fixed centres and frames.\par
\textbf{Authority/proof:} SA0 \S{}\S{}2--3; \S{}8.1; main text \S\ref{sec:attachment}.\par
\textbf{Prior ownership:} Yes: Paper B; SA0 types the inherited structure\par
\textbf{Qualification:} No continuous polygon field, deformation, seam/rim data or physical position follows.\par
\textbf{Interpretation limit:} Only the discrete section exists as a native state representation.\par
\textbf{Supporting predicates:} SA0 ordinals 176--211\par
\end{minipage}\par\medskip
\par\noindent\begin{minipage}{\linewidth}
\subsection*{L35 Two distinct seam criteria}
\textbf{EXACT THEOREM}. For constant face extensions, ambient equality at every seam requires \(q=0,p=be\); matched parallelism requires \(q=ae,p=be\).\par
\textbf{Domain:} Conditional extension criteria on the canonical three vertical planes.\par
\textbf{Authority/proof:} SA0 \S{}3.3; main text \S\ref{sec:attachment}.\par
\textbf{Prior ownership:} SA0 comparison, not new native seam law\par
\textbf{Qualification:} Neither condition is imposed; not a classification of every intrinsic notion of continuity.\par
\textbf{Interpretation limit:} No physical seam law is inferred.\par
\textbf{Supporting predicates:} SA0 ordinals 212--216\par
\end{minipage}\par\medskip
\par\noindent\begin{minipage}{\linewidth}
\subsection*{L36 Finite-carrier census}
\textbf{EXACT THEOREM}. Degree (0,1,2) counts are face scalar/normal (1,1,8), tangent (0,4,8), ambient-face (1,5,16), vertical seam (0,2,4), rim scalar (1,1,5), and rim vector (1,3,12).\par
\textbf{Domain:} Full D3h, polynomial state maps, finite carriers; seam labels via opposite face; whole-rim values.\par
\textbf{Authority/proof:} SA0 \S{}5.1--5.6; main text \S\ref{sec:attachment}.\par
\textbf{Prior ownership:} SA0; centre counts reused CM0\par
\textbf{Qualification:} Pointwise fields require a spatial function space; values on loops need not be tangent everywhere.\par
\textbf{Interpretation limit:} Multiplicity cannot supply physical selection.\par
\textbf{Supporting predicates:} SA0 ordinals 154--173\par
\end{minipage}\par\medskip
\par\noindent\begin{minipage}{\linewidth}
\subsection*{L37 Preparation type separation}
\textbf{DERIVED IDENTITY}. Lens and incident data pass through \(\C^2\otimes\C^3\), fixed \(U^n\), extraction \(\chi^\dagger\otimes I\), \(\Omega_0\), then \(D\Omega_0\). Preparation internals have no material shell attachment and no feedback enters F.\par
\textbf{Domain:} Named existing response, tensor order, branch and initialization receipt.\par
\textbf{Authority/proof:} SA0 \S{}7.1--7.2; main text \S\ref{sec:attachment}.\par
\textbf{Prior ownership:} Preparation formula prior Atlas 08; unified separation SA0\par
\textbf{Qualification:} Final Omega does have the existing D attachment; C\textasciicircum{}6 has twelve real dimensions.\par
\textbf{Interpretation limit:} Initialization is not boundary dynamics.\par
\textbf{Supporting predicates:} SA0 ordinals 223--226\par
\end{minipage}\par\medskip
\par\noindent\begin{minipage}{\linewidth}
\subsection*{L38 Fibre closure criterion}
\textbf{EXACT THEOREM}. \(\mathcal AF=F_{\mathcal A}\mathcal A\) on the image exists iff \(\mathcal A(x)=\mathcal A(x')\) implies \(\mathcal A(Fx)=\mathcal A(Fx')\).\par
\textbf{Domain:} Any set maps; restrict domain explicitly for local/regular versions.\par
\textbf{Authority/proof:} SA0 \S{}8.1; main text \S\ref{sec:closure}.\par
\textbf{Prior ownership:} Elementary quotient criterion reproved, no priority claim\par
\textbf{Qualification:} No smoothness, locality or invertibility conclusion.\par
\textbf{Interpretation limit:} Closure is not physical observability.\par
\textbf{Supporting predicates:} Analytic proof/source statement; no individually matched predicate claimed.\par
\end{minipage}\par\medskip
\par\noindent\begin{minipage}{\linewidth}
\subsection*{L39 Nonclosed lower readouts}
\textbf{EXACT THEOREM}. Single-face values, intensities, raw chirality, W, Gamma and joint (W,Gamma) generally fail closure. At \(g=0\), \(I_i^+=I_i[1+\eps(k_i-I_i)]^2\).\par
\textbf{Domain:} Exact native witnesses in SA0; parameters fixed in each comparison.\par
\textbf{Authority/proof:} SA0 \S{}8.1--8.2; main text \S\ref{sec:closure}.\par
\textbf{Prior ownership:} CM0 had ambient witnesses; these explicit fibre comparisons SA0\par
\textbf{Qualification:} Do not claim every combination/parameter locus is nonclosed.\par
\textbf{Interpretation limit:} Lost state information obstructs an autonomous readout law.\par
\textbf{Supporting predicates:} SA0 ordinals 219--222\par
\end{minipage}\par\medskip
\par\noindent\begin{minipage}{\linewidth}
\subsection*{L40 Complex coherence quotient}
\textbf{EXACT THEOREM}. \(\mathcal H=\Omega\Omega^\dagger\) closes when all diagonal entries of \(M\mathcal HM^T\) are positive: \(\mathcal H^+=U_\delta M\mathcal HM^TU_\delta^\dagger\), with delta from imaginary cubes of normalized transposed coherence entries.\par
\textbf{Domain:} Rank<=1 PSD Hermitian H; real native parameters; regular pre-sync domain.\par
\textbf{Authority/proof:} SA0 \S{}8.3, first theorem; main text \S\ref{sec:closure}.\par
\textbf{Prior ownership:} SA0 synthesis result\par
\textbf{Qualification:} One-step domain need not be invariant; zero-stratum obstruction retained.\par
\textbf{Interpretation limit:} Channel coherence is not a spatial metric.\par
\textbf{Supporting predicates:} Analytic proof/source statement; no individually matched predicate claimed.\par
\end{minipage}\par\medskip
\par\noindent\begin{minipage}{\linewidth}
\subsection*{L41 Real Gram quotient}
\textbf{EXACT THEOREM}. \(S_G=qq^T+pp^T\) also closes on the regular pre-stage domain: equal Gram factors differ by \(O(2)\), namely common phase and/or conjugation.\par
\textbf{Domain:} Rank<=2 real PSD Gram matrix; same regularity test.\par
\textbf{Authority/proof:} SA0 \S{}8.3, second theorem; main text \S\ref{sec:closure}.\par
\textbf{Prior ownership:} SA0\par
\textbf{Qualification:} S\_G loses chiral orientation; autonomy follows from symmetry, not injectivity.\par
\textbf{Interpretation limit:} A nonlocal channel quotient, not one scalar per face.\par
\textbf{Supporting predicates:} Analytic proof/source statement; no individually matched predicate claimed.\par
\end{minipage}\par\medskip
\par\noindent\begin{minipage}{\linewidth}
\subsection*{L42 Global zero-stratum failure}
\textbf{EXACT THEOREM}. For \(\eps=g=0,\lambda=1/10\), inputs \((1,-1,0)\) and \(e^{i\pi/6}(1,-1,0)\) have identical coherences and real Grams. Outputs have \(\mathcal H_{12}^+=-1\) versus \(-e^{-i/5}\), and \(S_{G,12}^+=-1\) versus \(-\cos(1/5)\).\par
\textbf{Domain:} Native Arg0 convention; same exact parameter tuple.\par
\textbf{Authority/proof:} SA0 \S{}8.4; main text \S\ref{sec:closure}.\par
\textbf{Prior ownership:} SA0 explicit closure witness; zero convention caveat prior Paper A/B\par
\textbf{Qualification:} At lambda=0 both Gram congruences close globally.\par
\textbf{Interpretation limit:} No repaired zero convention is introduced.\par
\textbf{Supporting predicates:} Analytic proof/source statement; no individually matched predicate claimed.\par
\end{minipage}\par\medskip
\par\noindent\begin{minipage}{\linewidth}
\subsection*{L43 Retained computational evidence}
\textbf{NUMERICAL EVIDENCE}. Counts are GR0 45, GR1 38, GR2 42, CM0 78, SA0 226, total 429 heterogeneous exact/source checks. Native finite differences, 100-digit evaluations and residuals remain separate datasets.\par
\textbf{Domain:} Only bounded fixtures and precisions actually recorded.\par
\textbf{Authority/proof:} ALL GR0 \S{}7; GR1 \S{}7; GR2 \S{}8; CM0 \S{}7; SA0 \S{}9; main text \S\ref{sec:repro}.\par
\textbf{Prior ownership:} Retained evidence, not newly rerun or new theorems\par
\textbf{Qualification:} Predicates are not proofs; small native residuals cannot resolve every tiny imaginary part.\par
\textbf{Interpretation limit:} No broad numerical survey or certification is claimed.\par
\textbf{Supporting predicates:} Analytic proof/source statement; no individually matched predicate claimed.\par
\end{minipage}\par\medskip
\par\noindent\begin{minipage}{\linewidth}
\subsection*{L44 Interpretation boundary}
\textbf{PHYSICAL INTERPRETATION}. Response inner products, screening, complex modes, circulation, vector parity, face sections and initialization do not by themselves define the corresponding physical metric, gravity, wave, current, position, deformation or flux.\par
\textbf{Domain:} Claims confined to the finite native mathematical model.\par
\textbf{Authority/proof:} ALL Checkpoint conclusions and explicit limits; main text \S\ref{sec:limits}.\par
\textbf{Prior ownership:} Consistent with accepted paper boundaries\par
\textbf{Qualification:} Absence in this model is not a universal no-go theorem.\par
\textbf{Interpretation limit:} No new physical theory is established.\par
\textbf{Supporting predicates:} Analytic proof/source statement; no individually matched predicate claimed.\par
\end{minipage}\par\medskip
\par\noindent\begin{minipage}{\linewidth}
\subsection*{L45 Open interfaces}
\textbf{OPEN QUESTION}. Global branch/zero-stratum analysis, nonlinear uniform limits, privileged observable/metric selection and minimal sufficient coherence subsets remain open.\par
\textbf{Domain:} No new research undertaken for this publication synthesis.\par
\textbf{Authority/proof:} ALL GR2 \S{}9; CM0 \S{}6; SA0 \S{}10 question 6; main text \S\ref{sec:discussion}.\par
\textbf{Prior ownership:} Open, not claimed results\par
\textbf{Qualification:} Existing finite proofs do not solve these interfaces.\par
\textbf{Interpretation limit:} Research publication stops at synthesis.\par
\textbf{Supporting predicates:} Analytic proof/source statement; no individually matched predicate claimed.\par
\end{minipage}\par\medskip
\endgroup


\section{Triad algebra}\label{app:triad}
The principal minors of \(M_3=\sigma R^{-1}-ee^T\) sum to
\[
\sum_{i<j}\left[\left(\frac{\sigma}{r_i}-1\right)
\left(\frac{\sigma}{r_j}-1\right)-1\right]
=\frac{\sigma(a^2+b^2+c^2)}{abc}=t.
\]
Its trace is h, and its zero eigenvalue is supplied by \(M_3r=0\). This derives \eqref{eq:triadpoly} without treating the sum-zero plane as invariant for unequal J. For \((a,a,c)\), the antisymmetric equal-pair mode has
\(\mu_A=2+c/a\), while the other transverse mode has
\(\mu_B=c/a+2a/c\); substituting in \eqref{eq:resonance} decides full-matrix collisions.

The alternating cycle expression and the imaginary Bloch coefficient are different invariants:
\[
\mathcal C_{\rm cyc}=\frac{g\ell dV_r}{a^2b^2c^2}\mathcal K,\qquad
\Gamma_B=\frac{\ell g[\ell(1+3g)-g]V_r}{abc}.
\]
In particular their exceptional parameter factors are not interchangeable. This paper uses the plus sign before \((3g\ell-g+\ell)\sigma_3\) in verified GR1 (20); the missing plus in the later GR2 (36) restatement is recorded as a transcription correction, not a new result or a modification of GR2.

\section{Bloch coefficient derivation}\label{app:bloch}
Set \(z\ne0\) formally before specializing to roots of unity. Compute J from \eqref{eq:bloch}, then
\[
T_B=\sum_iJ_{ii},\qquad
B_B=\sum_{i<j}(J_{ii}J_{jj}-J_{ij}J_{ji}),\qquad D_B=\det J.
\]
The first contains no z. The second is Laurent-linear in z:
\[
B_B(z)=B_B(1)+C_B(z+z^{-1}-2)-\frac{\Gamma_B}{2}(z-z^{-1}).
\]
The antisymmetric coefficient is obtained by subtracting the expression at \(z^{-1}\); its numerator is alternating in a,b,c and factors as
\(-\ell g[\ell(1+3g)-g]V_r(z-z^{-1})/(abc)\).
The symmetric coefficient reduces to the stated \(C_B\). The determinant follows from the product of the two stage determinants, proving \eqref{eq:BT}--\eqref{eq:BD}. These operations, rather than eigenvalue fitting, are reproduced by the paper-local checker.

At N=6 the blocks are \(1,-1\); at N=12 they are \(1,i,-1,-i\). Exact Fourier injection matrices have columns \(z^ne_\alpha\); direct multiplication \(J_NB_z=B_zJ(z)\) verifies the decomposition. The characteristic polynomial of the full matrix is the product of block polynomials, as this invariant direct sum spans every coordinate.

\section{Exact six-ring fixture}\label{app:fixture}
At \(\tau=1/1000\), \(\eps=1/20\), \(g=999999/4000000\), the native coefficients in \eqref{eq:fixture} are
\[
(k_1,k_2,k_3)=(491493/500000,\ 1,\ 508493/500000),\qquad
k=749993/750000.
\]
They lie strictly between \((49/50)^2\) and \((51/50)^2\).
Let \(L=49/50,U=51/50\). On the real box \(r_i\in[L,U]\), the amplitude map has off-diagonal derivatives g and diagonal derivatives
\(1+\eps(k_i-3r_i^2)-2g\), all nonnegative for this fixture and the interpolation of coefficients from their mean. Its row sum is
\(1+\eps(k_i-3r_i^2)\le9067893/10^7<1\).
Monotonicity reduces box invariance to the constant endpoints: their coupling vanishes, the lower endpoint increments are nonnegative because \(k_i>L^2\), and upper endpoint increments are nonpositive because \(k_i<U^2\). This proves contraction and existence/uniqueness in the box throughout the interpolation, connecting the exact fixture to the local positive equal branch.

At the defect, the row--column pairing in \eqref{eq:nilpotent} is
\((1+\tau)-2+(1-\tau)=0\). Since the matrix is nonzero, its minimal polynomial is \(\xi^2\) and characteristic polynomial \(\xi^3\). Exact rational substitution of \(\delta=\pm10^{-14}\) into \eqref{eq:unfoldcoeff} gives discriminants with the respective signs. Their full rational values are retained in GR2 and the package evidence; the sign is not inferred from rounded eigenvalues.

\section{Representation decomposition and observable kernels}\label{app:representations}
For an orthogonal finite-group representation X and target Y, averaging its action on \(\operatorname{Hom}(\operatorname{Sym}^dX,Y)\) projects onto equivariants. Its trace is the dimension:
\[
\frac1{12}\sum_{g\in D_{3h}}\chi_{\operatorname{Sym}^dX}(g)\chi_Y(g),\qquad
\chi_{\operatorname{Sym}^2X}(g)=\frac{\chi_X(g)^2+\chi_X(g^2)}2.
\]
This both proves the multiplicity computation and gives a finite exact verification route, with no numerical representation fitting.

For transverse u,v,
\[
\|N_f(u\odot u)\|^2=\frac14\|u\|^4,\quad
\|N_f(u\odot v)\|^2=\frac14\|u\|^2\|v\|^2,
\]
\[
\Gamma_{\rm ch}^2=3(\|u\|^2\|v\|^2-(u\cdot v)^2).
\]
These follow by \(u_3=-u_1-u_2,v_3=-v_1-v_2\) and the exact frame matrix. If \(N_fu=(X,Y,0)\), then
\[
N_f(u\odot u)=\tfrac13(2XY,X^2-Y^2,0);
\]
polarization gives \(N_f(u\odot v)=\tfrac13(XW+YZ,XZ-YW,0)\) for \(N_fv=(Z,W,0)\). The first is onto the plane as a complex-square map; the second is invertible in v for fixed nonzero u. Thus its kernel is u=0 or v=0. Gamma vanishes exactly for dependent transverse vectors. Mean-times-transverse maps vanish when either factor vanishes. Linear maps \(T_fu,be_z,T_fv,ae_z\) have the evident transverse/mean kernels. Images of mixtures cannot be obtained by independently adding images with shared inputs.

\section{Complete finite-carrier census}\label{app:census}
Define \(Kx=e\times x/\sqrt3\), \(K^2=-\Pi\), and
\[
\mathcal H_s=(a^2,b^2,\|u\|^2,\|v\|^2),\quad
\mathcal E_s=(au,\Pi(u^2),bv,\Pi(v^2)),\quad
\mathcal O_s=(av,bu,\Pi(uv)),
\]
with entrywise products in this appendix.

\paragraph{Face scalars and normal coefficients.}
Bases at degrees zero, one and two are \(e\), \(Ku\), and
\(\{he:h\in\mathcal H_s\}\cup\mathcal E_s\). Multiplying the i-th scalar by outward \(n_i\) gives the identical count for polar normal assignments.

\paragraph{Face tangent vectors.}
Represent outputs by coefficient pair (Q,P). Linear basis is
\((ae,0),(u,0),(0,be),(0,v)\). Quadratic basis is
\((Kf,0)\), \(f\in\mathcal E_s\), together with
\((0,\Gamma_{\rm ch}e)\) and \((0,Kh)\), \(h\in\mathcal O_s\).
There is no constant tangent map. Orthogonal normal/tangent decomposition gives counts \(1,5,16\) for ambient vectors per face.

\paragraph{Seams.}
Opposite-face indexing has the ordinary face permutation representation. Scalars and ambient vectors therefore have the corresponding face counts. Vertical seam coefficients transform like p and have linear basis \(be,v\) and quadratic basis
\(\Gamma_{\rm ch}e,K(av),K(bu),K\Pi(uv)\).

\paragraph{Whole rims.}
Let \(f_\sigma=f_e+\sigma f_o\), \(\sigma=\pm1\). Even scalars have bases 1 and \(\mathcal H_s\); odd scalars have b and \(\Gamma_{\rm ch}\), at degrees one and two. For polar vectors \(V_\sigma=V_e+\sigma V_o\), \(V_e\) has the main polar basis, while \(V_o\) has constant \(e_z\), linear \(N_fv\), and quadratic basis
\(\{T_fh:h\in\mathcal O_s\}\cup\{he_z:h\in\mathcal H_s\}\).
The horizontal reflection exchanges rims; under a general element \(Ge_z=\eta e_z\), the odd target transforms as \(\eta G\).

These lists are independent by mean/transverse and q-q/p-p/q-p separation; K is invertible on \(e^\perp\). The character counts above exhaust their dimensions. They require no selected seam or boundary point, but they do not supply a pointwise material field.

\section{Closure counterexamples and observable evolution}\label{app:closure}
The rational comparisons in Section~\ref{sec:closure} supply exact fibre failures, not only changes along one trajectory. The zero state is fixed and every homogeneous polynomial observable vanishes there. CM0's bounded examples can therefore be used in the same way. For example, at \(\eps=1/20,g=1/5,k=e,\lambda=0\),
\[
q=e,\ p=(1,-1,0)\quad\mapsto\quad
q^+=(19/20,19/20,1),\ p^+=(7/20,-7/20,0).
\]
Here \(T_fu\) and Gamma initially vanish, but the former is generated and
\(\Gamma_{\rm ch}^+=7/200\). At
\[
q=(1,-1,0),\quad p=(1,1,-2),
\]
one obtains \(q^+=(7/20,-7/20,0)\),
\(p^+=(7/20,7/20,-1/2)\). This generates the imaginary mean and W from zero, with \(W^+=(-7/200,7\sqrt3/200,0)\).
Swapping q and p supplies the corresponding complementary witnesses; a real \((1,2,3)\) maps to \((8/5,17/10,6/5)\), changing the direction of \(T_fq\).
The complete checkpoint list is retained; no trajectory scan is added.

For a homogeneous quadratic map \(V(x)=B(x,x)\), polarization gives
\(\Delta V=2B(x,\Delta x)+B(\Delta x,\Delta x)\). This is the exact common formula behind the individual increments. It supplies no independent observable law unless the fibre criterion is met.

\section{Verification inventory}\label{app:verification}
The complete measurements, including every stored high-precision digit and
all exact rational fixtures, are retained in \texttt{retained\_evidence.json}.
The following tables round values for print only. Maxima below compare like
quantities within a stated checkpoint, never across error types. The 429
retained exact/source predicates are enumerated separately; none is replaced
by these residuals. Paper-local checks are an additional bounded algebra and
reference audit, with their own machine-readable result.

\subsection*{Native complex128 finite differences}
For GR0 the error is the maximum entrywise error of the complete real
Jacobian. For GR1 it is separated into phase and radial blocks; for GR2 only
phase response was tested. These are different quantities.

\begin{center}\begin{tabular}{@{}lrr@{}}
\toprule
GR0 ring & Step & Complete Jacobian error\\\midrule
C3 & \(1.00000\times10^{-4}\) & \(5.00211\times10^{-10}\)\\
C3 & \(1.00000\times10^{-5}\) & \(1.33778\times10^{-11}\)\\
C3 & \(1.00000\times10^{-6}\) & \(6.12623\times10^{-11}\)\\
C12 & \(1.00000\times10^{-4}\) & \(5.02309\times10^{-10}\)\\
C12 & \(1.00000\times10^{-5}\) & \(2.87038\times10^{-11}\)\\
C12 & \(1.00000\times10^{-6}\) & \(1.62148\times10^{-10}\)\\
\bottomrule\end{tabular}\end{center}

\begin{center}\begin{tabular}{@{}lrrr@{}}
\toprule
GR1 N & Step & Phase error & Radial error\\\midrule
3 & \(1.00000\times10^{-4}\) & \(1.62181\times10^{-10}\) & \(5.00302\times10^{-10}\)\\
3 & \(1.00000\times10^{-5}\) & \(1.62181\times10^{-12}\) & \(7.91794\times10^{-12}\)\\
3 & \(1.00000\times10^{-6}\) & \(1.62093\times10^{-14}\) & \(1.63248\times10^{-10}\)\\
12 & \(1.00000\times10^{-4}\) & \(1.61339\times10^{-10}\) & \(5.00302\times10^{-10}\)\\
12 & \(1.00000\times10^{-5}\) & \(1.61340\times10^{-12}\) & \(5.35114\times10^{-12}\)\\
12 & \(1.00000\times10^{-6}\) & \(1.61537\times10^{-14}\) & \(1.63248\times10^{-10}\)\\
\bottomrule\end{tabular}\end{center}

\begin{center}\begin{tabular}{@{}lrr@{}}
\toprule
GR2 N & Step & Phase error\\\midrule
3 & \(1.00000\times10^{-6}\) & \(1.62093\times10^{-14}\)\\
6 & \(1.00000\times10^{-6}\) & \(1.61537\times10^{-14}\)\\
12 & \(1.00000\times10^{-6}\) & \(1.61537\times10^{-14}\)\\
\bottomrule\end{tabular}\end{center}

GR0 uses \(\eps=1/20,g=1/5,\lambda=1/1000\), with
\((k,\chi)=(1,0)\) on C3 and \((4,\pi/7)\) on C12. GR1/GR2 use the
weak positive branch at \(k=1,\kappa=(-1,0,1),\eta=.01\) and the same
parameters. Native fixed residuals are zero in GR0, and at most
\(1.11023\times10^{-16}\) in the two GR1 tests.
GR0's native frame/transport comparison and potential comparison each give
\(1.11023\times10^{-16}\), as two separate quantities.

\subsection*{100-digit transcription finite differences}
These checks use the mathematical transcription with mpmath, not a
100-digit execution of the complex128 native kernel.

\begin{center}\begin{tabular}{@{}lrr@{}}
\toprule
Method/domain & Step & Maximum error\\\midrule
GR0 C3 & \(1.00000\times10^{-10}\) & \(5.00000\times10^{-22}\)\\
GR0 C3 & \(1.00000\times10^{-20}\) & \(5.00000\times10^{-42}\)\\
GR0 C12 & \(1.00000\times10^{-10}\) & \(5.00000\times10^{-22}\)\\
GR0 C12 & \(1.00000\times10^{-20}\) & \(5.00000\times10^{-42}\)\\
GR1 phase & \(1.00000\times10^{-10}\) & \(1.62181\times10^{-22}\)\\
GR1 phase & \(1.00000\times10^{-20}\) & \(1.62181\times10^{-42}\)\\
\bottomrule\end{tabular}\end{center}

\subsection*{Weak-branch residuals and asymptotic coefficients}

\begin{center}\begin{tabular}{@{}lrr@{}}
\toprule
\(\eta\) & \(\|r-r_0-\eta u\|_\infty/\eta^2\) & \(\mathcal C_{\rm cyc}/\eta^3\)\\\midrule
0.01 & \(2.00516\times10^{-2}\) & \(-6.95490\times10^{-8}\)\\
0.005 & \(2.00477\times10^{-2}\) & \(-6.95491\times10^{-8}\)\\
0.0025 & \(2.00457\times10^{-2}\) & \(-6.95491\times10^{-8}\)\\
\bottomrule\end{tabular}\end{center}

GR1 fixed-equation residual has maximum \(1.00858\times10^{-99}\).

GR1 cycle-factorization error has maximum \(1.11591\times10^{-103}\).

GR1 pre-stage balance residual has maximum \(3.57184\times10^{-102}\).

The proved leading cycle coefficient is \(-6.95491\times10^{-8}\).

\begin{center}\begin{tabular}{@{}lrrr@{}}
\toprule
\(\eta\) & \(\mathcal D_3/\eta^2\) & \(\max|\im\rho|,\ z=i\) & \(\max|\im\rho|/\eta^3\)\\\midrule
0.01 & \(2.40509\times10^{-3}\) & \(4.21737\times10^{-13}\) & \(4.21737\times10^{-7}\)\\
0.005 & \(2.40510\times10^{-3}\) & \(5.27172\times10^{-14}\) & \(4.21738\times10^{-7}\)\\
0.0025 & \(2.40510\times10^{-3}\) & \(6.58966\times10^{-15}\) & \(4.21738\times10^{-7}\)\\
\bottomrule\end{tabular}\end{center}

GR2 fixed-equation residual has maximum \(1.01215\times10^{-99}\).

GR2 characteristic residual has maximum \(1.82317\times10^{-101}\).

The constructed triad SPD matrix has minimum eigenvalue \(9.98841\times10^{-1}\), balance residual \(1.60733\times10^{-101}\), and a simultaneous nonzero cycle obstruction \(-6.95490\times10^{-14}\). Its complete entries and all block eigenvalues are retained separately. These 100-digit quantities resolve the small imaginary parts; the native finite-difference residuals alone would not certify them.

\subsection*{Observable, quotient and preparation evidence}

\begin{center}\begin{tabular}{@{}lrr@{}}
\toprule
CM0 observable derivative & Step & Maximum over 12 observables\\\midrule
equal & \(1.00000\times10^{-5}\) & \(3.06800\times10^{-11}\)\\
weak unequal & \(1.00000\times10^{-5}\) & \(1.12197\times10^{-11}\)\\
\bottomrule\end{tabular}\end{center}

The separate CM0 partial-zero chirality formula error is \(1.78351\times10^{-17}\). The native decoder identity passed as a Boolean check, with no numerical residual invented.

\begin{center}\begin{tabular}{@{}p{.52\linewidth}rr@{}}
\toprule
SA0 identity & Residual & Tolerance\\\midrule
Complex coherence quotient identity & \(1.67250\times10^{-16}\) & \(2.00000\times10^{-13}\)\\
real Gram closure under phase & \(5.55112\times10^{-17}\) & \(2.00000\times10^{-13}\)\\
real Gram closure under conjugation & 0 & \(2.00000\times10^{-13}\)\\
zero-stratum witness formula & \(1.11022\times10^{-16}\) & \(2.00000\times10^{-13}\)\\
canonical FaceState step & 0 & \(1.00000\times10^{-15}\)\\
decoded view after step & 0 & \(1.00000\times10^{-15}\)\\
SRG handoff direct product n=0 & 0 & \(2.00000\times10^{-13}\)\\
SRG handoff direct product n=1 & \(1.09714\times10^{-17}\) & \(2.00000\times10^{-13}\)\\
SRG handoff direct product n=2 & \(3.12732\times10^{-18}\) & \(2.00000\times10^{-13}\)\\
SRG handoff direct product n=3 & \(6.60563\times10^{-18}\) & \(2.00000\times10^{-13}\)\\
\bottomrule\end{tabular}\end{center}

The SA0 zero-stratum real-Gram output gap is \(1.99334\times10^{-2}\). It is a nonclosure signal, not an approximation error.

\paragraph{Separation of evidence classes.}
Exact rational signs at the six-ring unfolding, integer symmetry dimensions,
and source-dependency predicates belong to the retained exact inventory.
Finite differences, nonlinear equation residuals, eigenvalue residuals,
constructed-metric residuals and quotient comparisons are numerical evidence.
The written proofs in the main text do not depend on interpreting one type
of measurement as another.

